\documentclass[../main.tex]{subfiles}
\begin{document}
% SOURCE p=166; printed=151; object=ch05-title; region=top
\chapter{泛函分析基础}\label{chapter:ch05-functional-analysis}
% SOURCE p=166; printed=151; object=ch05-introduction; region=middle
“泛函分析”是研究 $\mathbb R$ 或 $\mathbb C$ 上的无限维向量空间及其间线性映射的传统名称。它与纯粹线性代数的区别，在于拓扑方面的考虑具有重要地位。有限维向量空间上只有一种合理的拓扑，线性映射自动连续；但在无限维情形，事情就没有这么简单了。（正如已经指出的，若 $\{f_n\}$ 是 $\mathbb R$ 上的一列函数，陈述“$f_n\to f$”可以有许多不同含义。）本章只打算简要介绍这门学科，因此除\cref{section:ch05-topological-vector-spaces}外，我们只考虑由向量空间上的范数所定义的拓扑。
% SUPPLEMENT: id=v14-ch05-introduction-scope; parent=ch05-introduction; kind=motivation+hypotheses; origin=independent-derivation
\begin{remark}
\textbf{引言中的范围说明（重构新增）。}“合理的拓扑”需与向量运算相容并满足适当分离性；它不是说有限维集合上只有一个任意拓扑。本节\cref{exercise:ch05-finite-dimensional-norms}将完整证明所有有限维范数等价，并由此推出其间的线性映射连续。更一般的 Hausdorff 拓扑向量空间结论在\cref{section:ch05-topological-vector-spaces}讨论。无限维中则须明确选择哪种范数或拓扑，才有确定的收敛概念。
\end{remark}
% END SUPPLEMENT: id=v14-ch05-introduction-scope
\subfile{Chapter-05/section-01.tex}
\subfile{Chapter-05/section-02.tex}
\subfile{Chapter-05/section-03.tex}
\subfile{Chapter-05/section-04.tex}
\subfile{Chapter-05/section-05.tex}
\subfile{Chapter-05/section-06.tex}
\end{document}
